\justifying
\NSFCBodyText

\subsubsection{研究内容}
研究围绕“可用气象信息怎样改善概率预测，以及这种改善能否传递到爬坡预警”展开。拟以\textbf{【暂定15分钟采样、未来1--4小时预测】}作为短时任务，保持预测发布、输入获取和真实标签到达的时间关系。研究对象为单场站归一化功率，按正常并网运行和日照条件开展主体评价，另报全运行时段表现。日前预测仅作为取得业务数值天气预报（NWP）档案后的扩展，不与短时主任务并列承诺。

\textbf{研究内容一：多源气象信息与功率变化的跨尺度关系。}建立功率、观测、气象预报与场站运行状态的可追溯关联，分析不同气象变量在较慢变化与快速扰动中的作用。首先以功率历史和可计算的太阳几何信息形成最小输入，再逐类加入已可获得的地面观测、NWP及卫星信息，比较增益出现的天气状态与预测时距。对每类输入同时检查其测量或预报时刻与实际可取得时刻，从数据层排除未来信息。通过保持模型预算一致的变量和尺度消融，区分多源信息带来的改善与单纯增加容量的影响，形成用于后续模型选择的数据质量说明和尺度选择依据。

\textbf{研究内容二：天气变化下的概率功率预测与校准。}基于内容一保留下来的输入，建立有限规模的多尺度分位数预测模型，同时输出中位预测和一组功率分位数。按时序训练并以独立时间段选择模型，比较静态与滚动校准在晴稳、云量快速变化及有外部证据的沙尘过程中的表现。关注覆盖偏差、区间宽度和概率评分之间的关系，检查校准是否只改善多数天气，或通过过宽区间掩盖局部失配。形成逐时距的概率预测、校准诊断和历史完整预测误差向量，为爬坡概率计算提供边际与依赖两类输入。

\textbf{研究内容三：保持边际一致的爬坡概率估计与预警评价。}在各预测时距采用完全相同的分位数样本，仅改变跨时距样本的排序，构造独立排序、总体历史秩和天气条件历史秩三组轨迹。通过比较三组轨迹的功率净变化分布，分别检验时间依赖及天气条件选择对事件概率的贡献，并与直接预测增量的轻量方法比较。依据验证集确定的告警阈值，根据预测概率发出预警，分别评价概率质量和告警效用，最终形成可以回放天气过程、功率分布与预警结果的离线原型。

\subsubsection{研究目标}
总体目标是在真实信息可得性约束下，建立面向新疆典型光伏场站的短时概率预测与爬坡风险分析方法，并明确其适用天气和失效情形。形成三个可核验结果：一是说明哪些气象源和时间尺度相对功率历史基线提供稳定增益，证据来自同输入时限、同预算的跨时间段消融；二是量化滚动校准相对静态校准的覆盖偏差与区间效率变化，证据来自留出时段和天气分组的配对比较；三是检验时间依赖对爬坡事件概率及提前识别的贡献，证据来自同边际排序对照、直接增量基线和事件级不确定性分析。

项目以有无可重复的改善作为研究判断，不预先承诺固定精度提高百分比。即使某一复杂模型或依赖结构未优于简单对照，也将保留完整阴性结果和适用边界，避免只报告最有利场站、时距或天气。数据清单、模型和评价程序、案例报告及离线原型共同构成成果，地域化结论必须以新疆实测证据为前提。

\subsubsection{拟解决的关键科学问题}
第一个问题是：\textbf{在发布时刻可用的信息集合固定后，气象变量的预测增益怎样随时间尺度、天气状态和预报时距变化？}通过固定输入时限和计算预算比较尺度选择，判断增强是否来自可泛化信息。相关和滞后分析仅解释条件预测关系，不能替代物理因果识别。

第二个问题是：\textbf{天气转换造成误差分布变化时，利用已经到达的历史标签进行校准，能否兼顾覆盖水平和区间效率？}重点检查真实时间约束下的适应延迟与小样本问题。总体校准成功不等同于每类天气都获得覆盖保证，对少见天气将单独报告证据强度。

第三个问题是：\textbf{边际概率预测保持一致时，时间依赖对爬坡事件概率可靠性的贡献及其失效条件是什么？}若独立排序与保留依赖的轨迹具有相同逐时距分布而事件评分不同，就能把差异定位到时间依赖；若差异不稳定，则采用更简单的事件预测路线。这是本项目连接功率预测与预警的主要检验问题。

\subsubsection{研究方法}
\textbf{数据组织与时间切分。}拟获取\textbf{【暂定1--2个新疆场站、连续不少于12个月且尽量覆盖24个月的历史记录】}，优先包含交流侧功率、装机容量、辐照度、温度和运行状态；其他气象源按实际授权与可用性逐项增加。保存原始时间、统一时间及数据实际获取时刻，明确使用北京时间还是地方太阳时，避免新疆经度差异造成日照相位混淆。以装机容量归一化功率，保留容量调整记录。训练前核验功率单位、时间重复、跳变、缺测、检修和限电标记；无运行状态证据的异常不直接标为天气事件。

对时刻$t$的预测，只使用获取时刻不晚于$t$的数据。NWP按起报时刻、预报有效时刻和实际获取时刻匹配，卫星产品按实际处理延迟入模。ERA5再分析仅用于离线天气描述或明确标注的理想输入对照。缺少可追溯NWP时，主任务使用历史观测与功率，不把未来实况填入预测端。气象状态用于模型选择时须在$t$可知；未来实况天气标签只用于事后分层评价。夜间及停限机时段分别报告：事后获知的停机可以从气象机制分析中剔除，但应保留在全运行时段的预报评价中并说明原因。

采用连续时间块组成训练、验证、初始校准和测试集合，四者用途互不替代。数据条件允许时按时间先后以\textbf{【暂定50\%/15\%/15\%/20\%】}分配，并检查季节覆盖；不足两年时不声称完成跨年泛化验证。多时距样本仅在目标窗口完全落入所属集合时使用，边界隔离至少一个最大预测跨度；模型超参数和告警阈值在验证阶段确定，测试前锁定。条件允许时设置多个滚动起点及留站验证，单站结果不外推为所有新疆基地。归一化、插补参数、变量选择与小波尺度均只依赖训练信息；测试期间仅按预先确定的规则更新校准，模型选择不再使用测试效果。

\textbf{多尺度分位数建模。}以输入窗口内的单侧小波或因果滤波形成慢变与快变分量，任何边界处理均只使用$t$及以前的观测；不对包含测试未来值的全序列一次性分解。候选主模型在DWT-Former相关多尺度结构上设置分位数输出；WaveKAN作为有限备选，避免同时训练多个复杂模型。引入气象外生变量时比较固定融合与天气条件融合，保持可比的参数规模、训练轮数和输入长度。

设$p_{t+h}$为归一化实际功率，$\widehat q_{t,h}(\tau)$为$t$时刻对$t+h$的$\tau$分位预测，采用分位数损失
\[
\ell_\tau(y,q)=\bigl(\tau-\mathbf{1}\{y<q\}\bigr)(y-q).
\]
分位数网格采用\textbf{【暂定0.05--0.95、间隔0.05的分位水平】}，两端以已知物理界限作有界延伸，尾部估计另作敏感性分析。以单调输出或预先确定的分位数重排处理交叉，并将功率限制在当时已知容量对应的物理范围内。该处理并不证明预测分布正确，其效果仍需在独立时段检验。模型与持续性、晴空归一化持续性、线性预测、分位数梯度提升及一种标准时序模型比较，必要时减少深度模型数量，以保证对照和消融有完整重复记录。

\textbf{滚动校准与可靠性评价。}以目标覆盖率$1-\alpha$的原始区间为起点，对已获真实值的历史样本计算非一致性分数
\[
s_{i,h}=\max\{\widehat q_{i,h}(\alpha/2)-p_{i+h},\ p_{i+h}-\widehat q_{i,h}(1-\alpha/2)\}.
\]
在$t$时刻，只从满足$i+h\leq t$且真实值已经到达的样本中选取校准分数；缺测或迟报样本等待到齐后再用。用预定窗口的分位数修正区间端点，校准量取非负值以避免区间反转，再按固定规则约束到物理范围内，所有评分均评价处理后的最终区间。有限校准样本的分位位置采用共形化分位数回归的相应修正\cite{romano2019cqr}；逐时距更新，比较不校准、静态校准和滚动校准。校准窗口与步长先在历史验证中选择，\textbf{【暂定以80\%和90\%预测区间为主要报告层级】}。对时间相关和分布漂移条件，报告经验覆盖与区间评分，不承诺任意漂移下的有限样本覆盖保证。

逐时距计算平均分位数损失、覆盖偏差、归一化区间宽度和区间评分，并按天气与季节给出样本量。点预测另报归一化平均绝对误差（MAE）与均方根误差（RMSE），避免夜间零值使全日指标显得过好。分层样本不足时保留总体校准并单列该组不确定性，不事后调节分组直到指标合格。比较以同一测试时间块的配对差异为主，不把单个时间点当独立重复。

\textbf{同边际轨迹与爬坡定义。}设$\mathcal I_t$为$t$时刻可用信息，$\delta$为固定爬坡间隔，定义
\[
R_{t,h,\delta}=p_{t+h+\delta}-p_{t+h},\qquad
E^+_{t,h}=\mathbf{1}\{R_{t,h,\delta}>r\},\quad
E^-_{t,h}=\mathbf{1}\{R_{t,h,\delta}<-r\}.
\]
这里只评价两个端点之间的净变化，$h>0$且两个端点均落在预测采样网格上，$h+\delta$不超过最大预测跨度；不把它等同于窗口内最大瞬时变化。\textbf{【暂定以30分钟净变化超过装机容量20\%为主事件，10\%与30\%作敏感性分析】}，具体阈值在查看测试集前结合业务需求确认并固定。完成数据核验后、查看留出测试标签前，将主事件、告警选择规则、天气分组、模板回退和事件匹配参数保存为统一版本；尚未获得的业务认可不记为已确认。\textbf{【暂定以未来60至90分钟的下降爬坡为主要概率终点】}，其他时距和上升事件为次要或探索性检验。上下坡分别评价；跨夜间、缺测或停限机的标签按既定规则分层，防止把日出日落和非气象停机误解释为天气扰动。

轨迹构造采用已学习的单调分位数网格；两级区间校准是额外可靠性诊断，不能被视为整条累积分布已校准。对每个时距的同一预测分位数网格抽取同一组$M$个样本，保留依赖组根据历史完整误差向量的跨时距秩关系逐列排序，独立组则对各列独立置换。两组每一列的样本数值完全相同，只有行内时间配对改变，以隔离时间依赖贡献。历史模板借鉴Schaake shuffle的经验秩思想\cite{clark2004schaake,schefzik2013ecc}，所用残差来自历史滚动留出预测，不能采用模型在自身训练样本上的拟合残差。仅使用整条预测轨迹真实值均已到达的样本，将当前$t$时可知天气与各历史预测发布时已知的天气状态匹配，不按历史事后天气筛选模板；样本不足时使用总体模板。设置总体历史秩、发布时刻天气条件历史秩与独立随机排序三个对照，独立排序重复多个随机种子；\textbf{【暂定每次生成100条轨迹，并以50/200条检验抽样敏感性】}，模板不足时按事先固定的重采样规则处理并报告有效历史样本量。若总体历史依赖有效而天气条件选择无效，只接受前一层结论。根据第$m$条轨迹计算事件，得到
\[
\widehat\pi^+_{t,h}=\frac{1}{M}\sum_{m=1}^{M}
\mathbf{1}\{p^{(m)}_{t+h+\delta}-p^{(m)}_{t+h}>r\},
\]
下降概率同理。边际区间校准不自动保证联合事件概率校准，因此事件概率另用Brier评分、相对训练期天气/时段事件率基线的技能分数、可靠性图和事件样本量检验；路径关系辅以变差函数评分\cite{scheuerer2015variogram}。直接增量分位数或分类模型作为更简单的补充对照，并明确其没有“同边际”这一受控条件。

\textbf{预警评价与统计分析。}概率评分使用未经告警阈值化的全部有效概率。告警阈值只在验证集依据预设误报负担选择，测试集报告实际误报、精确率、召回率和提前量；不能根据测试结果重新调阈值来声称“固定误报”。把相邻重叠窗口聚合为同一事件段，按预先确定的时间容差进行一对一匹配，每起事件只计一次命中；提前量按首次有效告警相对事件起点计算，同时报告按窗口统计的结果，防止事件合并掩盖误报。稀有事件情况下比较精确率--召回率（PR）曲线与无技巧参考，避免以总体准确率评价。

采用连续日或天气过程为单位的配对分块自助法估计评分差异的不确定性，块长度根据训练期误差自相关选定并作敏感性检验。多时距和多天气结果完整报告，主要假设在测试前指定，探索性比较单列。主要检验固定为总体历史秩组减独立组的Brier差值$D_0$，天气条件历史秩组减总体历史秩组的差值$D_1$为预设次要检验，不根据测试表现交换二者顺序。各差值区间完全低于零才支持对应改善，跨零记为证据不足；若需判定实用等价，应预先另设容忍界限，不以不显著直接证伪。\textbf{【暂定每一重点天气或爬坡分层至少积累30个独立过程再作定量优劣判断】}；该数仅是报告门槛草案，不是统计功效保证，最终还需依据效应与区间宽度判断证据是否充分。

\subsubsection{技术路线}
\textbf{总体路线。}以获得合法使用条件的功率和气象数据为起点，先建立发布时刻的信息集合并识别可利用的时间尺度，内容一据此输出可追溯样本和变量选择依据。内容二在同一数据划分上训练概率模型并保存边际预测、到齐标签和连续误差向量；内容三接收这些产物，通过同边际轨迹对照检验爬坡概率，再将验证集确定的决策规则接入离线原型。三项内容在数据、概率输出和事件评价之间依次衔接，最终以独立时间段与天气过程检验总体目标。

\textbf{对应研究内容一的路线。}输入原始功率、气象与运行日志，完成授权核验、时间对齐和异常分层；以因果窗口提取尺度特征，在统一预算下进行逐源增量与变量消融。输出可复现数据清单、训练测试索引和尺度增益分析，用跨时间块的损失差异检验其稳定性。这些索引和保留变量成为内容二的固定输入。

\textbf{对应研究内容二的路线。}输入经核验的样本和尺度特征，训练分位数模型，并以只含到达标签的滚动样本校准区间。输出原始单调分位数网格、单独校准的区间、覆盖诊断和完整历史误差轨迹，比较静态与滚动校准的覆盖及效率。原始分位数网格与误差秩模板一并传给内容三，校准区间用于并行可靠性诊断；不将两个覆盖层级的端点拼接成已校准的完整分布，不把区间覆盖结果当作事件可靠性的替代证据。

\textbf{对应研究内容三的路线。}输入内容二的边际预测与可用历史误差轨迹，构造同边际、不同时间依赖的轨迹并计算爬坡概率，与独立排序和直接增量方法对照。输出事件概率、可靠性分析与按事件聚合的告警记录，以概率评分和锁定告警规则下的实际效果验证，最后形成可回放的离线原型和适用条件说明。

\subsubsection{关键技术}
第一项关键技术是记录气象信息何时可用，并限定尺度变换只使用过去数据，使模型输入与实际发布情景一致。每项变量保留取得时间，分解和插补仅使用过去信息，理想输入对照与实际可用输入明确分开。

第二项关键技术是处理真实值到达延迟的概率校准。不同预测时距对应不同可更新样本，只有标签到达后才更新对应分数；小样本天气以诊断为主，防止分组过细造成虚假可靠性。

第三项关键技术是同边际的依赖结构对照。逐列相同、逐行配对不同的轨迹使时间依赖贡献可以被单独检验，再通过独立事件统计连接预警效果。三项技术均以能够复现实验和发现失效条件为要求，不以模块数量作为评价依据。

\subsubsection{可行性分析}
本项目以申请人已确认的光伏多尺度预测和长时序研究为方法基础，利用已有计算资源开展有限规模实验。研究采用一个主预测模型、必要轻量对照和可复用后处理，不开展大模型预训练或新建传感网络。首期先核验数据并实现持续性及分位数基线，随后才决定复杂模型是否有足够证据进入比较，使两年任务与建议3万元的投入相匹配。

可行性的主要限制在数据和业务合作，而非模型训练算力。新疆实测功率和NWP档案尚待落实，相关缺口与应对措施在研究基础部分集中列示。若仅取得功率与历史观测，仍可开展短时预测和爬坡方法研究；如果始终缺少新疆实测，公开数据只能支持通用方法结论，项目的新疆验证目标应记为未完成，并据管理要求调整任务。此限制将在立项与中期检查时如实报告。
